\chapter{KIỂM THỬ VÀ ĐÁNH GIÁ HIỆU NĂNG}
\label{ch:chuong5}

\section{Phương pháp kiểm thử hệ thống}
\label{sec:phuong_phap_kiem_thu}

Để đánh giá tính đúng đắn, độ tin cậy và khả năng đáp ứng thực tế của hệ thống SmartDorm, quy trình kiểm thử toàn diện được thực hiện thông qua 4 phương pháp kiểm thử khoa học:
\begin{enumerate}
    \item \textbf{Kiểm thử đơn vị và tích hợp tự động (Unit \& Integration Testing):} Sử dụng framework Jest và thư viện Supertest trên nền tảng NestJS nhằm kiểm tra tính toàn vẹn của các hàm nghiệp vụ, logic băm thẻ RFID, quy tắc hoàn tác lệnh (Undo) và các điểm cuối API.
    \item \textbf{Đo đạc độ trễ truyền thông mạng (Network Latency Benchmark):} Đo lường thời gian trễ khứ hồi (RTT - Round-Trip Time) của các gói tin truyền qua giao thức MQTT Broker và WebSockets Socket.IO.
    \item \textbf{Kiểm thử khả năng chịu tải đồng thời (Concurrency \& Stress Testing):} Giả lập nhiều máy khách kết nối đồng thời để đánh giá mức độ ổn định và thời gian phản hồi của máy chủ backend khi có lượng lớn người dùng tương tác.
    \item \textbf{Kiểm thử phản ứng an toàn PCCC khẩn cấp (Emergency Fail-Safe Benchmark):} Đo đạc chính xác thời gian kích hoạt còi, thời gian tự động ngắt điện phòng và thời gian gửi thông báo Telegram khi nồng độ khí gas vượt ngưỡng nguy hiểm $\ge 80\%$.
\end{enumerate}

\section{Kết quả kiểm thử tự động toàn diện}
\label{sec:ket_qua_kiem_thu_tu_dong}

Hệ thống đã xây dựng và thực thi bộ kiểm thử hồi quy tự động gồm 84 ca kiểm thử (Test Cases) bao phủ toàn bộ các mô-đun chức năng. Bảng \ref{tab:tong_hop_kiem_thu} tổng hợp kết quả kiểm thử tự động của hệ thống.

\begin{table}[H]
\centering
\footnotesize
\begin{tabularx}{\textwidth}{|c|p{3.2cm}|c|c|c|X|}
\hline
\textbf{STT} & \textbf{Phân hệ kiểm thử} & \textbf{Số ca} & \textbf{Đạt} & \textbf{Lỗi} & \textbf{Nội dung kiểm tra cốt lõi} \\
\hline
1 & \texttt{AuthModule} & 12 & 12 & 0 & Xác thực JWT, mã hóa Bcrypt, bảo vệ route phân quyền RBAC \\
\hline
2 & \texttt{RoomsModule} & 10 & 10 & 0 & Quản lý danh sách phòng, liên kết tòa nhà và gán phần cứng ESP32 \\
\hline
3 & \texttt{DevicesModule} & 16 & 16 & 0 & Điều khiển 4 thiết bị, Command Pattern và cơ chế Hoàn tác (Undo) \\
\hline
4 & \texttt{AccessLogModule} & 14 & 14 & 0 & Thuật toán SHA-256 sinh mã thẻ RFID, đối soát cấp quyền ra vào \\
\hline
5 & \texttt{EmergencyModule} & 16 & 16 & 0 & Rule Engine PCCC, tự động ngắt điện khi gas $\ge 80\%$, xử lý sự cố \\
\hline
6 & \texttt{AnalyticsModule} & 8 & 8 & 0 & Tổng hợp dữ liệu vi khí hậu, tạo file xuất Excel và CSV chuẩn \\
\hline
7 & \texttt{Intelligence} & 8 & 8 & 0 & Mô hình Isolation Forest phát hiện bất thường và Trợ lý AI Copilot \\
\hline
\multicolumn{2}{|c|}{\textbf{TỔNG CỘNG}} & \textbf{84} & \textbf{84} & \textbf{0} & \textbf{Tỷ lệ đạt chuẩn: 100\% (84/84 ca Pass)} \\
\hline
\end{tabularx}
\caption{Bảng tổng hợp kết quả kiểm thử tự động các phân hệ chức năng của SmartDorm}
\label{tab:tong_hop_kiem_thu}
\end{table}

\section{Đánh giá độ trễ truyền thông mạng}
\label{sec:danh_gia_do_tre}

Độ trễ truyền nhận là chỉ số quyết định trực tiếp đến tính tức thời và độ an toàn của hệ thống IoT. Thử nghiệm được tiến hành qua 50 vòng đo liên tục giữa thiết bị phần cứng nhúng và máy chủ backend:

\begin{itemize}
    \item \textbf{Độ trễ truyền tin qua MQTT Broker (HiveMQ / Mosquitto):} Thời gian từ khi vi điều khiển ESP32 gọi lệnh \texttt{client.publish()} gửi gói tin telemetry đến khi máy chủ backend nhận được sự kiện qua kết nối TCP. Kết quả đo được: Độ trễ nhỏ nhất là $22.4\text{ ms}$, lớn nhất là $48.6\text{ ms}$, độ trễ trung bình đạt \textbf{34.2 ms} (thỏa mãn xuất sắc yêu cầu $< 50\text{ ms}$).
    \item \textbf{Độ trễ phát sóng WebSockets (Socket.IO):} Thời gian từ khi backend nhận dữ liệu đến khi giao diện trình duyệt web nhận được payload và cập nhật DOM. Kết quả đo được: Độ trễ trung bình đạt \textbf{11.8 ms} (thỏa mãn xuất sắc yêu cầu $< 15\text{ ms}$).
\end{itemize}

Bảng \ref{tab:do_tre_mang} tóm tắt các kết quả đo đạc độ trễ mạng trong điều kiện thử nghiệm thực tế.

\begin{table}[H]
\centering
\footnotesize
\begin{tabularx}{\textwidth}{|X|c|c|c|c|c|}
\hline
\textbf{Kênh truyền thông dữ liệu} & \textbf{Nhỏ nhất} & \textbf{Lớn nhất} & \textbf{Trung bình} & \textbf{Chỉ tiêu} & \textbf{Đánh giá} \\
\hline
Phần cứng ESP32 $\rightarrow$ MQTT Broker & 22.4 ms & 48.6 ms & \textbf{34.2 ms} & $< 50$ ms & \textbf{Đạt} \\
\hline
MQTT Broker $\rightarrow$ Backend NestJS & 4.1 ms & 8.5 ms & \textbf{5.6 ms} & $< 10$ ms & \textbf{Đạt} \\
\hline
Backend NestJS $\rightarrow$ Web Client (Socket.IO) & 7.2 ms & 14.8 ms & \textbf{11.8 ms} & $< 15$ ms & \textbf{Đạt} \\
\hline
Toàn trình (End-to-End từ Sensor lên Web) & 35.8 ms & 68.2 ms & \textbf{48.4 ms} & $< 100$ ms & \textbf{Đạt} \\
\hline
\end{tabularx}
\caption{Bảng đo đạc độ trễ truyền thông mạng trên hệ thống SmartDorm qua 50 vòng thử nghiệm}
\label{tab:do_tre_mang}
\end{table}

\section{Đánh giá khả năng chịu tải đồng thời}
\label{sec:danh_gia_chiu_tai}

Để bảo đảm hệ thống có thể mở rộng cho toàn bộ ký túc xá quy mô hàng ngàn sinh viên, thử nghiệm chịu tải đồng thời được thực hiện bằng cách giả lập các máy khách gửi truy vấn liên tục lên hệ thống:

\begin{table}[H]
\centering
\footnotesize
\begin{tabularx}{\textwidth}{|c|X|c|c|c|}
\hline
\textbf{Số kết nối đồng thời} & \textbf{Thời gian đáp ứng trung bình} & \textbf{Tỷ lệ thành công} & \textbf{Tải CPU máy chủ} & \textbf{Tải RAM} \\
\hline
50 kết nối & 24.5 ms & 100.0\% & 4.2\% & 145 MB \\
\hline
100 kết nối & 31.8 ms & 100.0\% & 6.8\% & 168 MB \\
\hline
250 kết nối & 48.2 ms & 100.0\% & 14.5\% & 210 MB \\
\hline
500 kết nối & 64.7 ms & 99.8\% & 26.1\% & 285 MB \\
\hline
\end{tabularx}
\caption{Kết quả kiểm thử khả năng chịu tải đồng thời của máy chủ backend SmartDorm}
\label{tab:kiem_thu_tai}
\end{table}

Kết quả cho thấy ngay cả khi có 500 kết nối đồng thời, thời gian phản hồi của hệ thống vẫn duy trì ở mức rất nhanh ($64.7\text{ ms}$), tỷ lệ thành công đạt 99.8\% và mức chiếm dụng tài nguyên CPU chỉ 26.1\%, khẳng định tính vững chắc của nền tảng phần mềm.

\section{Đánh giá độ tin cậy và thời gian phản ứng an toàn PCCC}
\label{sec:danh_gia_pccc}

Tính năng tự động ngắt điện phòng khi phát hiện rò rỉ khí gas $\ge 80\%$ là rào chắn an toàn quan trọng nhất của đề tài. Thử nghiệm thực nghiệm được tiến hành bằng cách kích hoạt khí gas nồng độ cao vào cảm biến MQ-2 và đo lường chính xác mốc thời gian diễn ra từng hành động bảo vệ:

\begin{table}[H]
\centering
\footnotesize
\begin{tabularx}{\textwidth}{|c|X|p{2.5cm}|p{2.5cm}|}
\hline
\textbf{Bước} & \textbf{Hành động bảo vệ của hệ thống} & \textbf{Thời gian phản ứng} & \textbf{Đánh giá} \\
\hline
1 & Cảm biến MQ-2 đạt ngưỡng nguy hiểm $\ge 80\%$ & $T_0 = 0.00\text{ s}$ & Điểm gốc đo \\
\hline
2 & ESP32 hú còi tại chỗ và \textbf{ngắt relay điện phòng} & $+0.085\text{ s}$ & Tức thời tại chỗ \\
\hline
3 & Backend tiếp nhận gói tin MQTT cảnh báo khẩn cấp & $+0.125\text{ s}$ & Mạng nội bộ \\
\hline
4 & Backend tạo sự cố \texttt{CRITICAL} \& phát sóng WebSockets & $+0.142\text{ s}$ & Web nháy đỏ \\
\hline
5 & Kích hoạt gửi tin nhắn Telegram Bot \& Email BQL & $+0.860\text{ s}$ & Nhận trên điện thoại \\
\hline
\multicolumn{2}{|l|}{\textbf{Thời gian ngắt điện phòng triệt tiêu tia lửa}} & \textbf{$< 0.1\text{ giây}$} & \textbf{Tuyệt đối an toàn} \\
\hline
\multicolumn{2}{|l|}{\textbf{Thời gian BQL nhận tin báo qua Telegram}} & \textbf{$< 1.1\text{ giây}$} & \textbf{Xuất sắc ($< 1.5$s)} \\
\hline
\end{tabularx}
\caption{Bảng đo đạc các mốc thời gian phản ứng của Cơ chế An toàn PCCC tự động}
\label{tab:thoi_gian_phan_ung_pccc}
\end{table}

Kết quả khẳng định cơ chế bảo vệ kép (Local Hardware Fail-Safe kết hợp Backend Emergency Broadcast) phản ứng cực kỳ nhanh nhạy: nguồn điện phòng được cắt đứt hoàn toàn chỉ sau $85\text{ ms}$, triệt tiêu triệt để nguy cơ phóng điện gây cháy nổ, và ban quản lý nhận được thông báo cứu nạn chỉ sau $0.86\text{ giây}$.

\newpage
